\documentclass[11pt]{article}

\usepackage[margin=1in]{geometry}
\usepackage{amsmath,amssymb}
\usepackage{enumitem}
\usepackage{hyperref}
\usepackage{xcolor}
\usepackage{booktabs}
\usepackage{listings}

\hypersetup{colorlinks=true, linkcolor=blue, urlcolor=blue, citecolor=blue}

\lstset{
  basicstyle=\ttfamily\small,
  breaklines=true,
  columns=fullflexible,
}

\title{TOAE209/TOAH209 -- Design and Analysis of Algorithms\\[4pt]
\large Week 5: Practical Exercises}
\author{Foreign Trade University}
\date{}

\begin{document}
\maketitle

\section*{Instructions}

For each problem below there is a starter file
\texttt{exercises/starter\_code\_problemNN.py} containing function signatures with
\texttt{raise NotImplementedError} bodies, and a small \texttt{\_\_main\_\_} block with
tests. Implement the missing functions so the tests pass. A reference solution is
provided in \texttt{solutions/solution\_problemNN.py} -- try the problem yourself before
looking!

All problems use only the Python 3.10+ standard library. Shared helper functions live in
\texttt{starter\_code.py} at the week root: \texttt{generate\_weighted\_graph},
\texttt{empty\_graph}, \texttt{add\_edge}, \texttt{edge\_list}, \texttt{vertices}, and the
\texttt{UnionFind} class (with path compression and union by rank). A graph is an
\emph{adjacency dict}: a map from each vertex to a dict \texttt{\{neighbour: weight\}};
undirected edges appear in both endpoints' dicts.

The 12 problems are grouped into three themes:
\begin{itemize}
    \item \textbf{Shortest Paths} (Problems 1--6): a weighted-graph builder, Dijkstra's
    algorithm (distances and path reconstruction), Bellman--Ford with negative edges,
    cross-validation, and a counter-example showing Dijkstra fails on a negative edge.
    \item \textbf{Minimum Spanning Trees} (Problems 7--11): Kruskal's algorithm with
    union-find, Prim's algorithm with a heap, verification that they agree, and the cut
    and cycle properties.
    \item \textbf{Applications} (Problem 12): single-link $k$-clustering via Kruskal.
\end{itemize}

\newpage
\section{Shortest Paths}

\subsection*{Problem 1 -- Weighted Graph Builder and Edge Listing}

Implement \texttt{build\_graph(edges, directed=False)} that builds an adjacency-dict graph
from a list of \texttt{(u, v, w)} edge tuples (using \texttt{add\_edge} from
\texttt{starter\_code.py}), \texttt{neighbors(graph, u)} returning \texttt{u}'s
neighbour$\to$weight dict (empty if isolated), \texttt{list\_edges(graph,
directed=False)} returning each edge once as \texttt{(u, v, w)}, and
\texttt{total\_weight(graph, directed=False)} summing all edge weights.

\textbf{Example.} \texttt{build\_graph([(0,1,4),(1,2,3),(0,2,5)])} yields a graph with
\texttt{g[0][1] == 4} and \texttt{g[1][0] == 4}; \texttt{total\_weight} is \texttt{12}.

\subsection*{Problem 2 -- Dijkstra's Shortest-Path Distances}

Implement \texttt{dijkstra(graph, source)} using a binary heap (\texttt{heapq}) to return
a dict of shortest-path distances from \texttt{source} to every vertex (non-negative
weights assumed; unreachable vertices get \texttt{float('inf')}).

\textbf{Hint.} Use the lazy heap version: push \texttt{(distance, vertex)} on each
improvement and skip a popped vertex that has already been finalized.

\subsection*{Problem 3 -- Dijkstra with Predecessors and Path Reconstruction}

Implement \texttt{dijkstra\_with\_pred(graph, source)} returning \texttt{(dist, pred)},
where \texttt{pred[v]} is the vertex preceding \texttt{v} on a shortest path (\texttt{None}
for the source and for unreachable vertices). Then implement
\texttt{reconstruct\_path(pred, source, target)} returning the vertex list
\texttt{[source, ..., target]}, or \texttt{None} if \texttt{target} is unreachable.

\subsection*{Problem 4 -- Bellman--Ford with Negative Edges}

Implement \texttt{bellman\_ford(graph, source, directed=False)}: relax all edges $n-1$
times, then do one extra pass to detect a negative cycle. Return the distance dict, or
\texttt{None} if a negative cycle is detected. Unlike Dijkstra, this handles negative
weights.

\textbf{Note.} For an undirected graph, relax each edge in both directions (and any
negative undirected edge is itself a negative cycle).

\subsection*{Problem 5 -- Verifying Dijkstra and Bellman--Ford Agree}

Implement \texttt{agree\_on\_random\_graph(n, seed)}: generate a random connected
non-negative graph (\texttt{generate\_weighted\_graph(n, seed=seed, edge\_prob=0.4,
max\_weight=15)}) and return \texttt{True} iff \texttt{dijkstra} (Problem 2) and
\texttt{bellman\_ford} (Problem 4) give equal distances from source 0. Then implement
\texttt{verify\_agreement(num\_trials, n, seed)} that returns \texttt{True} iff this holds
on \texttt{num\_trials} random instances (using \texttt{seed + trial}).

\subsection*{Problem 6 -- Counter-Example: Dijkstra Fails on a Negative Edge}

Implement \texttt{negative\_edge\_graph()} returning a small directed graph with one
negative edge (and no negative cycle) on which Dijkstra is wrong. Then implement
\texttt{dijkstra\_vs\_bellman\_ford()} returning \texttt{(dijkstra\_dist,
bellman\_ford\_dist)} for that graph; the two dicts must differ at some vertex.

\textbf{Hint.} Make Dijkstra \emph{finalize} a vertex at its direct-edge distance before
it ever explores the cheaper route that ends with the negative edge.

\newpage
\section{Minimum Spanning Trees}

\subsection*{Problem 7 -- Kruskal's MST with Union-Find}

Implement \texttt{kruskal(graph)}: sort the edges by weight and add each edge whose
endpoints lie in different components (tested with the \texttt{UnionFind} class). Return
\texttt{(mst\_edges, total\_weight)}.

\textbf{Example.} On a 4-vertex connected graph, the MST has exactly 3 edges; for the
graph in the starter tests its total weight is 6.

\subsection*{Problem 8 -- Prim's MST with a Heap}

Implement \texttt{prim(graph, start=None)}: grow a tree from \texttt{start} (default: the
first vertex), always adding the cheapest edge crossing from the tree to a vertex outside
it, using a binary heap. Return the total MST weight. Assume the graph is connected.

\textbf{Note.} The MST total weight is independent of the start vertex.

\subsection*{Problem 9 -- Verifying Kruskal and Prim Agree}

Implement \texttt{kruskal\_equals\_prim(n, seed)}: generate a random connected graph and
return \texttt{True} iff \texttt{kruskal} (Problem 7) and \texttt{prim} (Problem 8) give
the same total MST weight. Then implement \texttt{verify\_mst\_agreement(num\_trials, n,
seed)} returning \texttt{True} iff this holds on every random trial. (They may pick
different edge \emph{sets} when weights tie, but the total must match.)

\subsection*{Problem 10 -- Verifying the Cut Property}

Implement \texttt{crossing\_edges(graph, S)} returning all edges with exactly one endpoint
in the set \texttt{S}, \texttt{min\_crossing\_edge(graph, S)} returning a minimum-weight
crossing edge (or \texttt{None}), and \texttt{cut\_property\_holds(graph, S)} verifying
that the minimum crossing edge of the cut $(S, V\setminus S)$ belongs to some MST (build an
MST via Kruskal; in the tie case, force the edge in and confirm the total weight is
unchanged).

\subsection*{Problem 11 -- Classify an Edge: in EVERY / SOME / NO MST}

Implement \texttt{classify\_edge(graph, edge)} returning the string \texttt{"every"},
\texttt{"some"}, or \texttt{"none"} according to whether \texttt{edge} is in every, some
(but not all), or no MST. Compare the global MST weight with the MST weight when the edge
is \emph{forced in} and when it is \emph{forbidden}: forcing it worsens the weight
$\Rightarrow$ \texttt{"none"}; forbidding it worsens the weight $\Rightarrow$
\texttt{"every"}; otherwise \texttt{"some"}.

\textbf{Hint.} This is the computational form of the cut and cycle properties.

\newpage
\section{Applications}

\subsection*{Problem 12 -- Single-Link $k$-Clustering via Kruskal}

Implement \texttt{single\_link\_clustering(graph, k)}: run Kruskal's algorithm but stop
once exactly \texttt{k} components remain. Return \texttt{(clusters, spacing)}, where
\texttt{clusters} is a list of \texttt{k} lists of vertices and \texttt{spacing} is the
minimum inter-cluster edge weight (the next edge Kruskal would have added), or
\texttt{float('inf')} when \texttt{k} equals the number of vertices.

\textbf{Example.} For points with cheap edges $0$--$1$ and $2$--$3$ (weight 1) and an
expensive bridge $1$--$2$ (weight 10), the $2$-clustering is $\{0,1\}$ and $\{2,3\}$ with
spacing $10$.

\newpage
\section{LeetCode Practice}

After completing the problems above, practise this week's techniques (Dijkstra,
Bellman--Ford, and minimum spanning trees) on the following
\href{https://leetcode.com}{LeetCode} problems. Difficulty is LeetCode's own label.

\begin{itemize}
  \item \href{https://leetcode.com/problems/network-delay-time/}{Network Delay Time} -- \emph{Medium} -- Dijkstra.
  \item \href{https://leetcode.com/problems/path-with-minimum-effort/}{Path with Minimum Effort} -- \emph{Medium} -- Dijkstra / binary search.
  \item \href{https://leetcode.com/problems/cheapest-flights-within-k-stops/}{Cheapest Flights Within K Stops} -- \emph{Medium} -- Dijkstra / Bellman--Ford.
  \item \href{https://leetcode.com/problems/swim-in-rising-water/}{Swim in Rising Water} -- \emph{Hard} -- Dijkstra / union-find.
  \item \href{https://leetcode.com/problems/path-with-maximum-probability/}{Path With Maximum Probability} -- \emph{Medium} -- Dijkstra variant.
  \item \href{https://leetcode.com/problems/min-cost-to-connect-all-points/}{Min Cost to Connect All Points} -- \emph{Medium} -- MST (Prim / Kruskal).
  \item \href{https://leetcode.com/problems/connecting-cities-with-minimum-cost/}{Connecting Cities With Minimum Cost} -- \emph{Medium} -- MST (Kruskal).
  \item \href{https://leetcode.com/problems/find-the-city-with-the-smallest-number-of-neighbors-at-a-threshold-distance/}{Find the City With the Smallest Number of Neighbors at a Threshold Distance} -- \emph{Medium} -- shortest paths.
  \item \href{https://leetcode.com/problems/reconstruct-itinerary/}{Reconstruct Itinerary} -- \emph{Hard} -- greedy / Hierholzer.
\end{itemize}

\end{document}
